%% start of file `template.tex'.
%% Copyright 2006-2015 Xavier Danaux (xdanaux@gmail.com).
%
% This work may be distributed and/or modified under the
% conditions of the LaTeX Project Public License version 1.3c,
% available at http://www.latex-project.org/lppl/.


\documentclass[11pt,a4paper,sans]{moderncv}        % possible options include font size ('10pt', '11pt' and '12pt'), paper size ('a4paper', 'letterpaper', 'a5paper', 'legalpaper', 'executivepaper' and 'landscape') and font family ('sans' and 'roman')

  % moderncv themes
  \moderncvstyle{banking}                             % style options are 'casual' (default), 'classic', 'banking', 'oldstyle' and 'fancy'
  \moderncvcolor{blue}                               % color options 'black', 'blue' (default), 'burgundy', 'green', 'grey', 'orange', 'purple' and 'red'
  %\renewcommand{\familydefault}{\sfdefault}         % to set the default font; use '\sfdefault' for the default sans serif font, '\rmdefault' for the default roman one, or any tex font name
  %\nopagenumbers{}                                  % uncomment to suppress automatic page numbering for CVs longer than one page
  
  % character encoding
  %\usepackage[utf8]{inputenc}                       % if you are not using xelatex ou lualatex, replace by the encoding you are using
  %\usepackage{CJKutf8}                              % if you need to use CJK to typeset your resume in Chinese, Japanese or Korean
  
  % adjust the page margins
  \usepackage[scale=0.75]{geometry}

  % Add in acronym package and define them
  \usepackage[nolist,nohyperlinks]{acronym}
  \begin{acronym}
  \acro{ANAO}{Australian National Audit Office}
  \acro{HSF}{Herbert Smith Freehills}
  \acro{CPD}{Continuing Professional Development}
  \end{acronym}

  \usepackage{graphicx}
  \usepackage{wrapfig}
  \usepackage{xcolor}
  \graphicspath{ {./images/} }
  

  %\setlength{\hintscolumnwidth}{3cm}                % if you want to change the width of the column with the dates
  %\setlength{\makecvtitlenamewidth}{10cm}           % for the 'classic' style, if you want to force the width allocated to your name and avoid line breaks. be careful though, the length is normally calculated to avoid any overlap with your personal info; use this at your own typographical risks...
  
  % personal data
  \name{Mark}{Isherwood}
  % \title{Resumé}                               % optional, remove / comment the line if not wanted
  \address{Reservoir}{}{Victoria, Australia}% optional, remove / comment the line if not wanted; the "postcode city" and "country" arguments can be omitted or provided empty
  \phone[mobile]{+61 409 538 515}          
  \email{mark@markisherwood.me}                               % optional, remove / comment the line if not wanted
  \homepage{markisherwood.me}                         % optional, remove / comment the line if not wanted
  \social[linkedin]{mark-isherwood}                        % optional, remove / comment the line if not wanted
  % \social[github]{markisherwood}                              % optional, remove / comment the line if not wanted
  % \extrainfo{additional information}                 % optional, remove / comment the line if not wanted
  % \photo[64pt][0.4pt]{picture}                       % optional, remove / comment the line if not wanted; '64pt' is the height the picture must be resized to, 0.4pt is the thickness of the frame around it (put it to 0pt for no frame) and 'picture' is the name of the picture file
  % \quote{Some quote}                                 % optional, remove / comment the line if not wanted
  
  % bibliography adjustements (only useful if you make citations in your resume, or print a list of publications using BibTeX)
  %   to show numerical labels in the bibliography (default is to show no labels)
  \makeatletter\renewcommand*{\bibliographyitemlabel}{\@biblabel{\arabic{enumiv}}}\makeatother
  %   to redefine the bibliography heading string ("Publications")
  %\renewcommand{\refname}{Articles}
  
  % bibliography with mutiple entries
  %\usepackage{multibib}
  %\newcites{book,misc}{{Books},{Others}}
  %----------------------------------------------------------------------------------
  %            content
  %----------------------------------------------------------------------------------
\begin{document}
%\begin{CJK*}{UTF8}{gbsn}                          % to typeset your resume in Chinese using CJK
%-----       resume       ---------------------------------------------------------
\makecvtitle{}

\section{Experience}
  \subsection{Vocational}
    \cventry{2014--Present}{Senior Developer}{Wunderman-Bienalto}{Melbourne}{}{Melbourne Technology Team Lead.\newline{}%
    Detailed skills:
    \begin{itemize}
      \item Drupal:
        \begin{itemize}
          \item Worked on Drupal 7 and 8 websites.
          \item Experienced in custom module development
          \item Handled architecture for multiple Drupal websites
          \item Setup Apache Solr integration with Drupal and implemented custom sorting options
          \item Proficient with GovCMS
        \end{itemize}
      \item Adobe Target
        \begin{itemize}
          \item Developed countless experiences through Adobe Target for multiple clients
          \item Oversaw project to track users between multiple brands to better profile and target users for different experiences
        \end{itemize}
      \item Microsoft Dynamics:
      \begin{itemize}
        \item Setup multiple instances on Microsoft Dynamics 365 for different projects
        \item Created back-end services which worked with Dynamics through OData API
      \end{itemize}
      \item Build Systems
        \begin{itemize}
          \item Helped start and build a multi-purpose Gulp based build system for use on a wide variety of projects.
          \item Integrated a wide number of common build tasks (SASS/Babel/HTML tempting/Testing) into a single common project which could be used for almost any project
        \end{itemize}
      \item Accessibility
      \begin{itemize}
        \item Worked with several clients with a focus on accessibility requirements.
        \item Built and reviewed pages to ensure sites/pages met WCAG AA accessibility guidelines
      \end{itemize}
      \item Team Leadership
        \begin{itemize}
          \item Team Lead for the Melbourne development team
          \item Currently managing a team of 7 developers.
          \item Handling scheduling for the team and ensuring resources are properly utilised
          \item Mentoring newer developers
        \end{itemize}
      \end{itemize}}
    \cventry{2014--2014}{System Administrator}{cohealth}{Collingwood}{}{Assisted existing system administrator during large scale company merger\newline{}
    Detailed skills:
    \begin{itemize}
      \item Windows Server Administrator
      \begin{itemize}
        \item Worked maintaining an existing Windows based server environment for several health care offices
      \end{itemize}
      \item Thin Client Support
      \begin{itemize}
        \item Provided support for Citrix based thin clients across several office locations
      \end{itemize}
      \item Email and phone support
      \begin{itemize}
        \item Handled email and phone support for several different health care sites handling issues as required
      \end{itemize}
    \end{itemize}}
    \cventry{2012--2014}{Web Officer}{Swinburne University}{Glenferrie}{}{Started position as industry based learning year. Re-hired via contracting agency after inital placement.
    \begin{itemize}
      \item Custom PHP/MySQL Development
      \begin{itemize}
        \item Worked as sole maintainer of the online enrolment system for the university
      \end{itemize}
      \item Thin Client Support
      \begin{itemize}
        \item Provided support for Citrix based thin clients across several office locations
      \end{itemize}
      \item Email and phone support
      \begin{itemize}
        \item Handled email and phone support for several different health care sites handling issues as required
      \end{itemize}
    \end{itemize}}


\section{Education}
  \cventry{2011--2014}{Bachelor of Science (Professional Software Development)}{Swinburne University of Technology}{Glenferrie}{\textit{Distinction Average}}{}
  \cventry{2008--2010}{Diploma of Software Development and Website Development}{Swinburne University of Technology}{Glenferrie}{\textit{Distinction Average}}{}


\section{Certifications}
  \cvitemwithcomment{Sep 2015}{Acquia Certified Developer --- Drupal 7}{}
  \cvitemwithcomment{Sep 2017}{Microsoft Certified Professional}{}
  \cvitemwithcomment{Sep 2017}{Microsoft Dynamics CRM 2016 Customisation and Configuration}{}

\section{Computer skills}
  \cvitemwithcomment{Languages}{HTML, CSS, SASS, JavaScript (ES6), jQuery, C\#, Java, PHP, SQL}{}
  \cvitemwithcomment{Technologies}{AWS, Azure, Drupal, Node, Gulp, Yarn, Adobe Target, Apache Solr}{}

\section{References}
  \cvitemwithcomment{Alex Hobart}{Marketing Automation Lead}{+61 422 784 476}
  \cvitemwithcomment{Drew Cosgrove}{Former Senior Manager}{+61 411 393 354}
  Further references available upon request.

\clearpage
\section{Major Projects}
  \subsection{\href{https://www.engineersaustralia.org.au/}{Engineers Australia}}
    \begin{wrapfigure}{R}{0.125\textwidth}
      \vspace{-10pt}
      \includegraphics[width=0.125\textwidth]{EA_logo}
      \vspace{-20pt}
    \end{wrapfigure}
    Engineers Australia is a professional body for the engineering field. I worked to develop their online membership portal for their Drupal 7 site.
    
    The membership portal guided user's through the initial membership sign-up process, transferring to different membership tiers, renewal, and resignation. This included an administrative back-end for staff to process the applications through the workflow process.

    This project required extensive integration with Oracle based web services for handling business logic and data storage.

  \subsection{\href{https://www.anao.gov.au/}{\ac{ANAO}}}
    \begin{wrapfigure}{R}{0.125\textwidth}
      \vspace{-15pt}
      \colorbox{black}{\includegraphics[width=0.125\textwidth]{ANAO_logo}}
      \vspace{-25pt}
    \end{wrapfigure}
    The Australian National Audit Office's website was built using GovCMS\@. This project was conducted in close co-operation with Acquia. This provided a very interesting challenge on implementing all the required functionality with the limited modules approved for the distribution.

  \subsection{\href{https://www.herbertsmithfreehills.com/}{\ac{HSF}}}
    \begin{wrapfigure}{R}{0.2\textwidth}
      \vspace{-10pt}
      \includegraphics[width=0.2\textwidth]{HSF_logo}
      \vspace{-10pt}
    \end{wrapfigure}
    We developed a number of different websites for \ac{HSF} as part of a global re-platform. The two main sites we developed were Source and their corporate website.

    The Source website was built to allow legal professionals to track their yearly required professional development. The site hosted legal articles and videos which earned users \ac{CPD} points. A custom module was developed to provide integration with a third party video service which included tracking how much of a video had been watched before crediting points.

    The Corporate website was the largest of all the sites built. This involved a re-build of their existing corporate website (hosted on Sitecore) with additional functionality and content migration. One of the core features of this site is the search functionality which is featured across the entire website. This was built with a Apache Solr (hosted by Acquia Search) back-end to ensure all the required functionality was available and quick. This included search term auto-complete, spelling suggestions, acronym/synonym support, inline results, and faceted search.

    The corporate website also offered multi-lingual support for content allowing staff to have their profiles and articles in their native languages.


% \clearpage
%-----       letter       ---------------------------------------------------------
% recipient data
% \recipient{Company Recruitment team}{Company, Inc.\\123 somestreet\\some city}
% \date{January 01, 1984}
% \opening{Dear Sir or Madam,}
% \closing{Yours faithfully,}
% \enclosure[Attached]{curriculum vit\ae{}}          % use an optional argument to use a string other than "Enclosure", or redefine \enclname
% \makelettertitle

% Lorem ipsum dolor sit amet, consectetur adipiscing elit. Duis ullamcorper neque sit amet lectus facilisis sed luctus nisl iaculis. Vivamus at neque arcu, sed tempor quam. Curabitur pharetra tincidunt tincidunt. Morbi volutpat feugiat mauris, quis tempor neque vehicula volutpat. Duis tristique justo vel massa fermentum accumsan. Mauris ante elit, feugiat vestibulum tempor eget, eleifend ac ipsum. Donec scelerisque lobortis ipsum eu vestibulum. Pellentesque vel massa at felis accumsan rhoncus.

% Suspendisse commodo, massa eu congue tincidunt, elit mauris pellentesque orci, cursus tempor odio nisl euismod augue. Aliquam adipiscing nibh ut odio sodales et pulvinar tortor laoreet. Mauris a accumsan ligula. Class aptent taciti sociosqu ad litora torquent per conubia nostra, per inceptos himenaeos. Suspendisse vulputate sem vehicula ipsum varius nec tempus dui dapibus. Phasellus et est urna, ut auctor erat. Sed tincidunt odio id odio aliquam mattis. Donec sapien nulla, feugiat eget adipiscing sit amet, lacinia ut dolor. Phasellus tincidunt, leo a fringilla consectetur, felis diam aliquam urna, vitae aliquet lectus orci nec velit. Vivamus dapibus varius blandit.

% Duis sit amet magna ante, at sodales diam. Aenean consectetur porta risus et sagittis. Ut interdum, enim varius pellentesque tincidunt, magna libero sodales tortor, ut fermentum nunc metus a ante. Vivamus odio leo, tincidunt eu luctus ut, sollicitudin sit amet metus. Nunc sed orci lectus. Ut sodales magna sed velit volutpat sit amet pulvinar diam venenatis.

% Albert Einstein discovered that $e=mc^2$ in 1905.

% \[ e=\lim_{n \to \infty} \left(1+\frac{1}{n}\right)^n \]

% \makeletterclosing

%\clearpage\end{CJK*}                              % if you are typesetting your resume in Chinese using CJK; the \clearpage is required for fancyhdr to work correctly with CJK, though it kills the page numbering by making \lastpage undefined
\end{document}


%% end of file `template.tex'.
